\section{Removing completion in coheight one}\label{h3:sec:coheight-one}

The calculation of the completed height-two stratum must be compared with
the algebraic coefficient algebra occurring in ordinary Morava cooperations.
This section proves that comparison.  The two main points are finiteness at
each complete local-field stage and the use of the \emph{full} reduced-norm
quotient.  In particular, averaging over central tame units fails at
$p=7$.

Put
\[
 K=k((y)),\qquad J=P_2,\qquad
 H=\ker(\hthreeNrd:P_3\longrightarrow\Z_p^\times),\qquad
 H'=\hthreeNrd^{-1}(\mu_{p-1}).
\]
Let $L$ be the connected-marking algebra of
\Cref{h3:prop:integral-frames,h3:lem:marking-descent}.  Its commuting $P_3$- and
$J$-actions make it an ind-\'etale $J$-torsor over $K$.
For a normal open subgroup $U\subset J$, write $B_U=L^U$.
It is a finite Galois \'etale $K$-algebra with group $J/U$; in particular,
it is a finite product of complete discretely valued fields.  All field
valuations are normalized by $v(y)=1$.  Statements about valuation ideals
in a product are imposed on every factor.  Both $P_3$ and $J$ preserve
these valuations, allowing permutations of the factors.

We write $C(T,M)=C_{\hthreects}(T,M)$ for continuous functions from a
profinite space $T$ to a topological module $M$, and use this notation
termwise for cochain complexes. We also put
\[
 R=L^{\mathrm{perf}}=\colim_{U,e}B_U^{1/p^e},\qquad
 B_U^b=\widehat{B_U^{\mathrm{perf}}}.
\]
As throughout the paper, a cochain complex with coefficients in an
algebraic union means the filtered colimit of the cochain complexes of
the displayed complete stages.  Thus, for example,
\begin{equation}\label{h3:co:filtered-convention}
 C^\bullet_{\hthreects}(H,R)
 :=\colim_{U,e}C^\bullet_{\hthreects}(H,B_U^{1/p^e}).
\end{equation}
We do not identify this complex with continuous functions into a union
endowed with its subspace metric.  Group cohomology in negative degrees
is understood to be zero.

The input from \Cref{h3:thm:completed-cohomology} has the following precise
form.  Each $H^a_{\hthreects}(H,B_U^b)$ is finite.  Moreover,
\begin{equation}\label{h3:co:completed-input}
 H^a\!\left(\colim_U C^\bullet_{\hthreects}(H,B_U^b)\right)
 =
 \begin{cases}
 k,&a=0,\\
 k(2\epsilon),&a=3,\\
 0,&a\ne0,3,
 \end{cases}
 \qquad \epsilon\in\{1,-1\}.
\end{equation}
The degree-zero map is the constants map.  In the notation $k(2\epsilon)$,
an element of the full quotient $H'/H=\mu_{p-1}$ represented by a norm
unit $u$ acts as $u^{2\epsilon}$.  The comparison, the commuting $J$-action,
and the filtered transition maps are retained with every profinite
parameter space.  No statement about Galois descent of a single completed
infinite field is required below.

\begin{theorem}[Coheight-one comparison]\label{h3:thm:coheight-one-comparison}
For every $p\geq5$, the natural constants map is a quasi-isomorphism
\begin{equation}\label{h3:co:main-constants}
 k\xrightarrow{\ \sim\ }
 \colim_U C^\bullet_{\hthreects}(H',B_U).
\end{equation}
It is equivariant for the commuting connected-frame action and for the
remaining height-three action.  For every profinite space $T$, its
parameterized version
\[
 C(T,k)\xrightarrow{\ \sim\ }
 \colim_U C\bigl(T,C^\bullet_{\hthreects}(H',B_U)\bigr)
\]
is also a quasi-isomorphism.  All maps are the maps induced by constants
and inclusions of coefficient algebras.
\end{theorem}

The proof has three stages.  First, the distribution algebra of the
connected Honda marking torsor and its Cartier residue transpose prove
finiteness at every sufficiently deep complete local-field stage.
Second, Frobenius contracts the positive valuation lattices and compares
the algebraic perfection with its completion.  Finally, the full norm
quotient removes the remaining nonconstant row and the distribution
target.  Every comparison in the conclusion is the original constants
map, even though the proof uses auxiliary operators.

\subsection{Continuous cohomology and finite resolutions}

We first establish the topological facts needed for both the present
comparison and the ordinary-stratum argument of \Cref{h3:sec:ordinary}.
A linearly topologized $\F_p$-vector space below is complete and separated;
an action has invariant neighborhoods if its open linear subspaces have
a cofinal subsystem preserved by the group.

\begin{lemma}[Finite resolutions and continuous duality]
\label{h3:co:compact-resolution}
Let $G$ be a compact $p$-adic analytic group with no element of order
$p$, and put $d=\dim G$.  The trivial module $\F_p$ has a resolution of
length $d$ by finitely generated projective $\F_p[[G]]$-modules.
For every complete linearly topologized representation $M$ with invariant
neighborhoods, applying continuous $\F_p[[G]]$-linear Hom to this
resolution computes $C^\bullet_{\hthreects}(G,M)$ up to a continuous chain
homotopy equivalence.  Its terms are continuous direct summands of finite
powers of $M$.

If $M$ is compact, $H^a_{\hthreects}(G,M)$ is compact Hausdorff.
For $G=H$ or $G=P_3$, put $d_H=8$ and $d_{P_3}=9$.  Their
adjoint orientation characters are trivial, and there are natural
identifications
\begin{equation}\label{h3:co:compact-duality}
 H^a_{\hthreects}(G,M)^\vee
 \cong H^{d_G-a}_{\hthreects}(G,M^\vee),
 \qquad
 M^\vee=\Hom_{\hthreects,\F_p}(M,\F_p),
\end{equation}
where $M^\vee$ has its discrete topology and contragredient action.
The transposed coefficient map induces the dual cohomology map.
The same formula, with compact and discrete sides interchanged, holds
for a discrete continuous $\F_p$-representation and its compact dual.
In particular, finite coefficient modules have finite cohomology.
\end{lemma}

\begin{proof}
The groups needed here satisfy the hypotheses by
\Cref{h3:lem:framework-groups}; their open subgroups inherit the absence
of $p$-torsion and have the same Lie dimension.
For the general assertion, use the Noetherian finite-global-dimension
theorem for the completed
group ring \cite[Proposition~3.3]{AW} and the fact that its integral
dualizing cohomology is a free
rank-one module in degree $d$ and vanishes in the other degrees, with
the adjoint orientation action
\cite[Proposition~4.16 and Remark~4.23]{BGHS}.
These statements give a finite projective resolution over
$\Z_p[[G]]$.  Noetherianity makes its terms finitely generated.  Give each
term its compact topology as a summand of a finite power of the group
ring.  The maps and their images are continuous and closed.  Its syzygies
are $p$-torsion-free, so reduction modulo $p$ is exact.
The integral completed bar resolution is projective in compact modules:
the completed free $\Z_p$-module on a profinite set is compact and
torsion-free, hence projective, since its Pontryagin dual is divisible
and therefore injective.  Induction to $\Z_p[[G]]$ preserves that
projectivity.  The continuous projective comparison therefore shows
that the dual of the finite integral resolution computes the quoted
continuous dualizing cohomology.  Its reduction modulo $p$ does so as
well, because both its terms and its unique integral cohomology module
are $p$-torsion-free.  The reduced dual complex
has cohomology only in degree $d$.  If the reduced resolution extends
beyond $d$, its last dual differential surjects onto a projective module;
splitting and dualizing removes a contractible pair.  Repetition gives
the stated length-$d$ resolution.

Here is why the resolution computes the asserted topological cochains.
For a profinite space $Z$, continuous maps $Z\to M$ identify with
continuous linear maps $\F_p[[Z]]\to M$.  Modulo an invariant open
subspace of $M$, the map has finite image and its linear extension
factors through a finite clopen partition of $Z$.  Taking the inverse
limit over these subspaces gives the claimed extension into $M$.
The completed bar resolution is projective in compact modules: a
surjection of compact $\F_p$-vector spaces admits a continuous linear
section, obtained by dualizing and splitting the corresponding injection
of discrete vector spaces.  Such a section lifts the profinite set of
bar generators; the lift extends group-ring linearly.  The projective
comparison construction therefore gives continuous chain maps and
homotopies between the completed bar resolution and the finite
resolution.  Applying continuous Hom into $M$ proves the assertion.
Continuity for the uniform cochain topologies can be checked after
quotienting by each invariant open subspace.

For compact coefficients, the finite complex has compact terms and
closed boundaries, which proves the compact Hausdorff assertion.
Now take $G=H$ or $G=P_3$.  The adjoint orientation of each is trivial
by \Cref{h3:lem:framework-groups}: conjugation has determinant one on
$D_3$, and also on its trace-zero summand.  Thus the dualizing module
has the trivial right action, in degree $d_G=8$ or $9$ respectively.

Let $Q_\bullet$ be the finite resolution for this $G$, and set
$Q_i^*=\Hom_{\F_p[[G]]}(Q_i,\F_p[[G]])$.
The reversed group-ring dual resolves the trivial right module, by
the dualizing calculation above.  Regard $Q_i^*$ as a left module
through the involution $g\mapsto g^{-1}$.  Finite-projective evaluation
gives
\[
 \Hom_{\hthreects,\F_p[[G]]}(Q_i,M)^\vee
 \cong
 \Hom_{\hthreects,\F_p[[G]]}(Q_i^*,M^\vee).
\]
For a finite free module this is evaluation coordinate by coordinate;
passing to a projective summand preserves the identification.  It
commutes with both differentials and coefficient maps.  Continuous
$\F_p$-duality is exact on compact vector spaces, as follows by extending
a functional on a finite-dimensional quotient.  Reversing degrees about
$d_G$ and taking cohomology proves \eqref{h3:co:compact-duality} and its
naturality for both groups.  For a discrete continuous representation,
each element has a finite $G$-orbit; finite sets therefore lie in finite
$G$-stable submodules.  Their annihilators give invariant open
neighborhoods in the compact dual.  Dualizing the compact formula again
proves the discrete version.  Finally, the finite-resolution complex
has finite terms when $M$ is finite.
\end{proof}

Here the coefficient actions have the required invariant neighborhoods.
For the local-field and root stages, valuation balls are invariant.
For the ordinary spaces $B_{\hthreehol}$ and $B_{\hthreepk}$ of
\Cref{h3:sec:ordinary}, write $A=k[[x,y]]$ and use the valuations
$v_{a,b}$ with $v_{a,b}(x)=a$ and $v_{a,b}(y)=b$.  Since
$g(x)=xu_g$ with $u_g\in A^\times$ and $g(y)\in(x,y)$,
\begin{equation}\label{h3:co:uniform-action-bound}
 v_{a,b}(gf)\geq v_{a,\min(a,b)}(f)
 \qquad(g\in P_3,\ a,b>0).
\end{equation}
Indeed units and their inverses have valuation zero, so the estimate
holds termwise on Laurent polynomials.  Density extends it to
$B_{\hthreehol}$, then to its finite root levels and $B_{\hthreepk}$.
The coefficient lines used in \Cref{h3:sec:ordinary} have unit action
multipliers in their free $A$-lattices, so the same estimate holds
in those frames.  Given a basic open linear neighborhood $V$,
\eqref{h3:co:uniform-action-bound} supplies a neighborhood contained in
$\bigcap_{g\in P_3}g^{-1}V$.  This intersection is therefore open,
invariant, and contained in $V$.  The bound also holds uniformly for
continuous functions on a compact cochain domain, with any profinite
parameter space.  Thus these spaces satisfy the hypothesis of
\Cref{h3:co:compact-resolution} for the actual uniform cochain topologies.

An exact coefficient sequence induces an exact sequence of ordinary
continuous cochain complexes whenever its surjection has a continuous
section as a map of spaces.  Equivariance and additivity of the section
are unnecessary: compose a cochain with the section.  We will use this
for discrete valuation quotients and for finite-dimensional linear maps
over complete local fields, including their Frobenius transports.  In a
filtered union it suffices that every target stage have a continuous
lift at one larger source stage, and that each stage's kernel lie in a
kernel stage with its subspace topology.

\begin{lemma}[Strictness]\label{h3:co:strictness}
Let $A^\bullet$ be a complex of Polish abelian groups with continuous
differentials.  If $H^a(A^\bullet)$ is finite, its quotient topology is
discrete.  More precisely, the boundary subgroup $B^a$ is open and closed
in the cycle group $Z^a$, and
$\partial:A^{a-1}\twoheadrightarrow B^a$ is open.  Consequently, for
every neighborhood $V$ of zero in $A^{a-1}$ there is a neighborhood $W$
of zero in $Z^a$ such that every member of $W$ has a primitive in $V$.
\end{lemma}

\begin{proof}
The cycles form a closed Polish group.  The boundary subgroup is an
analytic subset of it, being the continuous image of a Polish space,
and hence has the Baire property \cite[Theorem~21.6]{Kechris}.
Its finite index implies that it is nonmeagre.  A nonmeagre set with the
Baire property is comeagre in some nonempty open set; intersecting two
sufficiently small translates of that open set shows that its difference
set contains a neighborhood of zero.  Applied to the subgroup $B^a$,
this proves openness, and an open subgroup is also closed.

We also recall the open-mapping step.  If $q:E\twoheadrightarrow F$ is
a continuous surjective homomorphism of Polish groups and $V$ is a
neighborhood of zero in $E$, choose open $V'$ with $V'-V'\subset V$.
Separability gives a countable covering of $E$ by translates of $V'$.
Thus $q(V')$ is an analytic nonmeagre subset of $F$, and its difference
set contains a neighborhood of zero.  This difference set lies in
$q(V)$.  Hence $q$ is open.  Apply this to the differential with target
$B^a$, which is already open in $Z^a$.
\end{proof}

\begin{lemma}[Profinite parameters]\label{h3:co:parameters}
Let $T$ be a profinite space.
\begin{enumerate}
\item A continuous map $T\to F$ lifts through every continuous open
surjective homomorphism $E\twoheadrightarrow F$ of Polish abelian groups.
\item If a Polish complex $A^\bullet$ has finite cohomology in every
degree, then
\[
 H^a C(T,A^\bullet)\cong C(T,H^a(A^\bullet)),
\]
where the finite cohomology group on the right is discrete.
\item For a filtered system of such complexes, the corresponding
identity commutes with the cochain colimit, with the colimit of the
cohomology groups given its discrete topology.
\end{enumerate}
\end{lemma}

\begin{proof}
For the first assertion, choose complete translation-invariant metrics
on the groups and successively smaller neighborhoods of zero in $E$
whose diameters tend to zero and whose sums converge.  Openness provides
corresponding neighborhoods in $F$.  Compactness and zero-dimensionality
give a finite clopen partition on which the given map is approximated
by finitely many chosen images.  Lift those values to $E$.
Repeat on the error, refining the partition and choosing the correction
inside the next prescribed neighborhood of zero in $E$.  The successive
locally constant lifts converge uniformly to a continuous lift; their
images converge to the original map.  This construction uses no
metrizability assumption on $T$.

A continuous family of cycles gives a continuous family of cohomology
classes.  If that family is zero, the first assertion and
\Cref{h3:co:strictness} lift it continuously through
$A^{a-1}\twoheadrightarrow B^a$.  Conversely, a continuous map from
$T$ to the finite discrete cohomology group has finite image; choose a
cycle for each value to obtain a continuous family.  This proves the
second assertion.  For the third, a continuous function to a discrete
filtered colimit has finite image.  Its finitely many values and equality
relations are represented at one stage.  Filtered colimits of vector
spaces are exact, which proves the claimed compatibility.
\end{proof}

All coefficient spaces to which \Cref{h3:co:strictness} is applied are
Polish.  A finite product of local fields with finite residue fields is
Polish.  Its perfected completion has a countable dense subset obtained
from the finite root levels.  For a compact metrizable profinite space
$Z$, $C(Z,M)$ is Polish when $M$ is Polish: uniform limits give
completeness, and finite clopen partitions with a countable dense subset
of $M$ give separability.  The groups $H^a$ serving as cochain domains
are compact metrizable.  The parameterized complex $C(T,A^\bullet)$
itself is not asserted to be Polish for arbitrary $T$.

We shall also use the following elementary countability fact.  If $G$ is
compact and second-countable and $N$ is countable discrete, then every
$C^a_{\hthreects}(G,N)$ is countable.  Indeed $G^a$ has countably many clopen
sets, and a cochain is specified by a finite clopen partition and
finitely many values.  For discrete coefficients the same finite-image
description shows that continuous cochains commute with filtered
colimits, even when the transition maps are not injective.

\subsection{Distribution operators and the invariant differential}

The normalized generator-lift torsors of
\Cref{h3:lem:marking-descent} are torsors over $L$
under the constant group schemes $\Gamma_2[p^l]$.  Their coordinate
algebras and transition embeddings are
\begin{equation}\label{h3:co:root-torsors}
 R_l=L^{1/p^{2l}},\qquad R_l\longrightarrow R_{l+1}.
\end{equation}
Here the transition on the torsors is multiplication by $p$.
The group $H$ commutes with translations, and $J$ conjugates them by
continuous constant automorphisms.  These assertions concern the finite
flat torsors, not just their geometric points.  In particular, their
radicial rank $p^{2l}$ is retained under the separable connected-marking
extension.  Before perfection, \Cref{h3:geom:regular-levels} gives the
regular finite algebras of this rank, and \Cref{h3:lem:marking-descent}
identifies their generator-lift algebra after the separable base change.

\begin{lemma}[Distribution parameter]\label{h3:co:distribution}
After a finite enlargement of $k$, the inverse-limit distribution
algebra of \eqref{h3:co:root-torsors} is $k[[D]]$, with a parameter $D$
having the following properties.
\begin{enumerate}
\item It acts $L$-linearly and locally nilpotently on $R$.  For every
$i\geq0$,
\[
 \ker(D^{p^i}:R\to R)=L^{1/p^i}.
\]
It is surjective, and $D^{p^i-1}:L^{1/p^i}\to L$ is nonzero on every
factor.
\item It commutes with $H$.  Its residual $\mu_{p-1}$-character is
$\chi(u)=u^{s_{\mathrm{dist}}}$, where $s_{\mathrm{dist}}\in\{1,-1\}$.  Thus
\begin{equation}\label{h3:co:distribution-sequence}
 0\longrightarrow L\longrightarrow R
 \xrightarrow{D} R(\chi^{-1})\longrightarrow0
\end{equation}
is an exact sequence of $H'$-modules.
\item If $q=p^i$, where $i$ is even and divisible by $[k:\F_p]$, then
\begin{equation}\label{h3:co:frobenius-distribution}
 D(f^q)=(D^qf)^q,\qquad
 D^q(z^{1/q})=(Dz)^{1/q}.
\end{equation}
In particular $D^q$ is $L^{1/q}$-linear.
\item On any complete finite separable and root stage these operators
are continuous after passage to a larger such stage.  Each target stage
admits a continuous additive right inverse of $D$ with values in one
larger stage.  Therefore \eqref{h3:co:distribution-sequence} is exact on
the filtered continuous cochain complexes, with profinite parameters.
\end{enumerate}
\end{lemma}

\begin{proof}
The Cartier dual of the height-two, dimension-one Honda group is again
connected of height two and dimension one. Indeed, on the Honda model
$\Phi^2=[p]$, so duality gives $V^2=[p]$ on its dual. Together with
$\Phi V=[p]$, cancellation on the divisible group gives $\Phi=V$;
thus each finite dual torsion kernel is killed by a Frobenius iterate
and is connected. The dimension formula and formal-group construction
in \cite[Section~2.2, Proposition~1, and Section~2.3, Proposition~3]{Tate1967PDivisible}
give dimension one and a power-series coordinate, already over $\F_p$.
No identification with a chosen coordinate law on the original Honda
group is needed. A parameter on the dual identifies
the finite distribution algebra with
\[
 \mathcal O(\Gamma_2[p^l])^*
 =k[D]/(D^{p^{2l}}).
\]
The duals of the torsor transition maps are the inclusions of the dual
torsion kernels.  Their restriction maps give the inverse limit
$k[[D]]$.

After a faithfully flat trivialization of a torsor, its functions form
the dual regular module of this truncated polynomial ring.  The pairing
given by the coefficient of $D^{p^{2l}-1}$ in a product is perfect, so
that module is one regular nilpotent block.  Ranks descend under faithful
flatness.  Consequently the kernel of $D^e$ on $R_l$ has rank $e$ over
$L$ for $1\leq e\leq p^{2l}$.

For $q=p^i$, the formal-group coproduct satisfies
\[
 \Delta(D^q)\in(D^q\otimes1,1\otimes D^q).
\]
The module-algebra identity therefore makes $\ker D^q$ a subalgebra.
On a field factor it is an intermediate field of degree $q$.
The element $y$ is a $p$-basis of $K$ and remains a $p$-basis after
separable extension: if $E/K$ is finite separable, then
$[E:E^p]=[K:K^p]=p$, and $y\notin E^p$ because a separable extension
contains no nontrivial purely inseparable subextension of $K$.
An intermediate field of degree $q$ in $E^{1/p^{2l}}/E$ lies in
$E^{1/q}$, since every element has purely inseparable degree at most
$q$; equality follows from the degrees.  Applying this on every factor
and then passing to $L$ proves the kernel formula.  The block calculation
also proves nonvanishing of the top operator.  In a larger block the
old kernel is contained in the image of $D$, which proves surjectivity
on $R$.

By \Cref{h3:lem:marking-descent}, a full reduced-norm unit $u$ rescales the
normalized \'etale generator by $u$ or $u^{-1}$.  On the fixed connected
framing cover this conjugates translations by the corresponding scalar
automorphism.  Cartier duality preserves scalar multiplication, so the
dual tangent character is $\chi(u)=u^{s_{\mathrm{dist}}}$ with $s_{\mathrm{dist}}=\pm1$.
Since $H$ acts trivially on translations, this is an action of the
quotient $H'/H$, whether or not that quotient has a subgroup section.
Starting with a parameter $D_0$, replace it by
\[
 D=\frac1{p-1}\sum_{u\in\mu_{p-1}}
       \chi(u)^{-1}u(D_0).
\]
Its linear term is the linear term of $D_0$, and it is an eigenparameter.
This proves the second assertion, including the inverse twist on the
target of \eqref{h3:co:distribution-sequence}.

For the Frobenius assertion, the Honda model over $\F_p$ satisfies
$\Phi^2=[p]$.  On the full divisible group,
$V\Phi=[p]=\Phi^2$; cancellation of the faithfully flat epimorphism
$\Phi$ gives $V=\Phi$.  Cartier duality exchanges the two operators.
For the stated even $i$, their $i$th iterates are $[p^{i/2}]$.
To make the scalar and duality conventions explicit, put
$A_l=\mathcal O(\Gamma_2[p^l])$ and pair it with its finite dual by
$\langle b,\xi\rangle=\xi(b)$.  Since $q$ fixes $k$, Cartier-dual
naturality for the $i$th relative Frobenius reads
\begin{equation}\label{h3:co:frobenius-dual-pairing}
 \langle b^q,D\rangle
   =\langle b,(V^i)^*D\rangle
   =\langle b,D^q\rangle.
\end{equation}
For the second equality, choose an $\F_p$-coordinate $t$ on the Cartier dual,
which descends to $\F_p$.  Tame averaging leaves $D=\sum_n c_nt^n$ defined over
$k$, and hence
$(V^i)^*D=\sum_n c_nt^{nq}=D^q$ because $c_n^q=c_n$.
Here multiplication in the distribution algebra corresponds to
composition of operators.  If a finite torsor coaction is written
$\rho(f)=\sum_j f_j\otimes b_j$, \eqref{h3:co:frobenius-dual-pairing}
therefore gives
\[
 D(f^q)=\sum_j f_j^qD(b_j^q)
       =\sum_j f_j^qD^q(b_j)
       =\left(\sum_j f_jD^q(b_j)\right)^q.
\]
The last equality uses $q$-power invariance of the constants.
The coefficients $f_j$ retain their $q$th powers throughout; this finite
Hopf-coaction calculation requires no splitting of the torsor.
This proves \eqref{h3:co:frobenius-distribution}.  If $a\in L^{1/q}$,
apply it to $af$ and use $a^q\in L$ and $L$-linearity of $D$ to obtain
$D^q(af)=aD^qf$.  The cancellation above was on a divisible group;
no cancellation on a finite torsion kernel has been used.

On a fixed root level only finitely many powers of $D$ act.  In the
basis of successive fractional powers of the $p$-basis $y$, their
matrices have finitely many coefficients in $L$, and these lie in one
finite separable stage.  The resulting matrices between finite-dimensional
spaces over complete local fields are continuous.  To construct a
right inverse on a stage, choose a $D$-preimage of each member of its
finite field basis and put those preimages in one larger complete
separable/root stage.  Linear extension gives a continuous additive
section.  The kernel at a finite stage has the required subspace
topology; intersection with the separable union is that stage's
separable part, by uniqueness of purely inseparable roots.  The
cochain-exactness criterion preceding \Cref{h3:co:strictness} proves the
last assertion.  All of these arguments are factorwise for a finite
product, with finitely many choices at each stage.
\end{proof}

The injection of $L$ into the filtered completed coefficient algebra,
together with degree zero of \eqref{h3:co:completed-input}, gives
\begin{equation}\label{h3:co:fixed-constants}
 L^H=k.
\end{equation}
Indeed an invariant element belongs to an $H$-stable finite stage;
completion injects on its invariants, and filtered degree-zero cohomology
of the target is exactly $k$.

The $p$-basis calculation also identifies ordinary and continuous
differentials: on a factor $\kappa((s))$ both are freely generated by
$ds$, since derivations kill $p$th powers and the constant field is
perfect.  Thus Cartier below is the intrinsic field operator, with
its usual continuous Laurent-series formula.

Let $\eta$ be the nowhere vanishing differential of
\Cref{h3:lem:invariant-differential}.  It is defined at a finite separable
stage and is $H$-invariant.  The canonical-line and conormal
identifications used there give, on this slice,
\[
 \Omega^1_{K/k}\cong\omega_K^{\otimes(p+2)}[\det].
\]
The connected marking trivializes $\omega_L$, and the full norm kernel
fixes the determinant factor.  We write $\mathrm{Car}$ for Cartier on
differentials and set
\begin{equation}\label{h3:co:cartier-definition}
 \mathcal C(f)=\frac{\mathrm{Car}(f\eta)}{\eta}.
\end{equation}
This is the coheight-one Cartier operator.  It is distinct from the
ordinary root projection $\mathcal P$ of \Cref{h3:sec:ordinary}.

\begin{lemma}[Cartier and distribution traces]\label{h3:co:cartier-trace}
For every $i\geq1$, put $q=p^i$.  There exists $c_i\in k^\times$ such that
\begin{equation}\label{h3:co:cartier-distribution-trace}
 D^{q-1}(f^{1/q})=c_i\mathcal C^i(f)\qquad(f\in L).
\end{equation}
In particular,
\begin{equation}\label{h3:co:cartier-frobenius-zero}
 \mathcal C(f^p)=0.
\end{equation}
At every finite separable stage $B$ containing $\eta$ there is an exact
sequence of continuous $\F_p[H]$-modules
\begin{equation}\label{h3:co:cartier-four-term}
 0\longrightarrow B\xrightarrow{\mathrm{Fr}}B
 \xrightarrow{d/\eta}B\xrightarrow{\mathcal C}B\longrightarrow0.
\end{equation}
Here and throughout this section, \(\mathrm{Fr}(f)=f^p\) is coefficient
\(p\)-Frobenius.  The ordinary-stratum argument in
\Cref{h3:sec:ordinary} separately fixes a \(q\)-power operator.
The two short exact factors of \eqref{h3:co:cartier-four-term} admit
continuous additive sections onto their images.  Under the residue pairing
\begin{equation}\label{h3:co:residue-pairing}
 \langle f,g\rangle_B
 =\sum_j\hthreeTr_{\kappa_j/\F_p}\hthreeRes_{s_j}(fg\eta),
 \qquad B=\prod_j\kappa_j((s_j)),
\end{equation}
the $H$-action is isometric and
$\langle\mathcal C f,g\rangle_B=\langle f,g^p\rangle_B$.
\end{lemma}

\begin{proof}
The functional $I_i(f)=D^{q-1}(f^{1/q})$ takes values in $L$, is nonzero
on every factor, is $H$-equivariant, and satisfies
$I_i(a^qf)=aI_i(f)$.  The same semilinearity and equivariance hold for
$\mathcal C^i$.  The Cartier pairing gives an isomorphism
\begin{equation}\label{h3:co:perfect-cartier-pairing}
 \mathrm{Fr}^i_*L\longrightarrow
 \Hom_L(\mathrm{Fr}^i_*L,L),\qquad
 b\longmapsto\bigl[f\longmapsto\mathcal C^i(bf)\bigr].
\end{equation}
Here $a$ acts on $\mathrm{Fr}^i_*L$ as multiplication by $a^q$.
For a field factor, both sides have dimension $q$; the map is injective,
since for $b\ne0$ one can choose $z$ with $\mathcal C^i(z)\ne0$ and
evaluate at $b^{-1}z$.  This proves the isomorphism on products and
separable filtered unions as well; equivalently it is the perfect
coefficient-extraction pairing in the common $p$-basis.
Thus there is a unique multiplier $b_i$ with
$I_i(f)=\mathcal C^i(b_if)$.  It is a unit because both functionals
are nonzero on every field factor.  Equivariance and uniqueness give
$b_i\in L^H=k$, by \eqref{h3:co:fixed-constants}.  Semilinearity gives
\eqref{h3:co:cartier-distribution-trace}, with $c_i=b_i^{1/q}$.

This comparison applies in particular when $i=1$; it does not require
the even-iterate identity \eqref{h3:co:frobenius-distribution}.  Since $D$
kills $L$, it yields
$c_1\mathcal C(f^p)=D^{p-1}f=0$, proving
\eqref{h3:co:cartier-frobenius-zero}.  The latter is therefore a consequence
of the invariant torsor and differential, not an identity for an
arbitrary normalization of Cartier.

On a factor $\kappa((s))$ the logarithmic Laurent formula is
\begin{equation}\label{h3:co:cartier-laurent}
 \mathrm{Car}\!\left(\sum_n a_ns^n\frac{ds}{s}\right)
 =\sum_n a_{pn}^{1/p}s^n\frac{ds}{s}.
\end{equation}
It proves continuity and surjectivity.  The kernel consists of the forms
whose coefficients in indices divisible by $p$ vanish; a continuous
additive primitive is $\sum_{p\nmid n}(a_n/n)s^n$.
The kernel of $d$ is $B^p$, and its root map is continuous.
A continuous section of Cartier sends
$\sum_n a_ns^n ds/s$ to $\sum_n a_n^ps^{pn}ds/s$.
Multiplication or division by the fixed nonzero $\eta$ transports these
sections to \eqref{h3:co:cartier-four-term}.  The maps themselves are
$H$-equivariant because $\eta$ is invariant and Cartier is intrinsic.

Residue is invariant under change of uniformizer.  The sum of finite
field traces is preserved by residue-field automorphisms and factor
permutations.  This proves $H$-invariance of
\eqref{h3:co:residue-pairing}.  Finally,
$\hthreeRes\mathrm{Car}(\alpha)=(\hthreeRes\alpha)^{1/p}$ and
$g\mathrm{Car}(\alpha)=\mathrm{Car}(g^p\alpha)$.
The finite-field trace is unchanged by $p$th roots.  Applying these
identities to $\alpha=f\eta$ proves the transposition formula.
\end{proof}

\begin{lemma}[Good finite stages]\label{h3:co:good-stages}
There is a cofinal family of normal open subgroups $U\subset J$ such
that, for $B=B_U$, the differential $\eta$ is defined over $B$, there is
$\theta\in B^{1/p}$ with $D\theta=1$, and
\begin{equation}\label{h3:co:good-jet}
 \sigma(D)-D\in D^{p+2}k[[D]]\qquad(\sigma\in U).
\end{equation}
The operator $D$ restricts continuously to $B^{1/p}$.
Writing $\eta=h_j(s_j)ds_j/s_j$ on each field factor, the lattice
\begin{equation}\label{h3:co:positive-differential-lattice}
 P_B=\prod_j h_j^{-1}s_j\kappa_j[[s_j]]
\end{equation}
is compact, open, and preserved by $H$ and $\mathcal C$.
Its annihilator for \eqref{h3:co:residue-pairing} is $\mathcal O_B$, and
the annihilator of $\mathcal O_B$ is $P_B$.
\end{lemma}

\begin{proof}
The nilpotent block on $L^{1/p}$ puts $1$ in the image of $D$, so choose
$\theta$ there.  The finite data $\theta^p$ and $\eta$ lie at a finite
separable stage.  The action of $J$ on the finite set of jets
$k[[D]]/(D^{p+2})$ is continuous.  Intersect their open stabilizers with
any prescribed open subgroup and pass to a normal core.  This proves
cofinality and \eqref{h3:co:good-jet}.  For $f\in B^{1/p}$, $D^pf=0$ and
\[
 \sigma(Df)-Df=(\sigma(D)-D)f=0.
\]
The kernel filtration and uniqueness of roots therefore put $Df$ in
$(L^{1/p})^U=B^{1/p}$.  Its finite matrix over the complete stage proves
continuity.

The condition defining $P_B$ says that $f\eta$ is an integral ordinary
differential on each factor.  That condition is independent of the
uniformizer and is preserved by $H$.  Formula
\eqref{h3:co:cartier-laurent} preserves strictly positive logarithmic
order, proving Cartier stability.  The annihilator statements follow
coefficient by coefficient: an integral $g$ pairs trivially with all
strictly positive logarithmic coefficients, while every pole of $g$
is detected by a suitable positive coefficient of $f\eta$.  Conversely
every nonpositive coefficient of $f\eta$ is detected by a monomial in
$\mathcal O_B$.  Nondegeneracy of the finite-field trace makes these
detections nonzero when needed.
\end{proof}

\subsection{Finiteness at the local-field stages}

We now use the distribution and Cartier operators to prove the missing
finite-stage bound.  The only finiteness assumed at this point is for
the completed perfected fields.  Compact residue duality first controls
the stable Cartier image of a lattice; a separate Frobenius-kernel
argument then forces the uncompleted field cohomology to be finite.

An order lattice in a finite product of local fields means a product of
fractional powers of their maximal ideals.

\begin{proposition}[Finite-stage finiteness]\label{h3:prop:finite-stage-finiteness}
Let $B=B_U$ be a good stage as in \Cref{h3:co:good-stages}.
For every $a\geq0$, each of the following cohomology groups is finite:
\[
 H^a_{\hthreects}(H,B),\qquad
 H^a_{\hthreects}(H,\Lambda),\qquad
 H^a_{\hthreects}(H,B/\Lambda),
\]
where $\Lambda$ is any $H$-stable order lattice in $B$.
The assertion holds at each finite root stage and for a coefficient line
which has an $H$-invariant frame at a finite separable stage, with its
order lattices and their quotients.  In particular it holds for the
periodicity and determinant lines used in this paper after choosing a
sufficiently deep stage.
\end{proposition}

We prove the proposition without assuming finite cohomology of the
uncompleted fields.

\begin{lemma}[Positive completed coefficients]\label{h3:co:completed-positive}
Let $B$ be a good stage.  All $H^a_{\hthreects}(H,B^b)$ are finite, and
\[
 R\Gamma_{\hthreects}(H,(B^b)^{++})=0,
 \qquad (B^b)^{++}=\{f:v(f)>0\}.
\]
The cohomology groups of $\mathcal O_{B^b}$ and
$B^b/\mathcal O_{B^b}$ are finite as well.
\end{lemma}

\begin{proof}
The first assertion is the finite-stage statement of
\Cref{h3:thm:completed-cohomology}.  Let $z$ be a continuous positive-valued
$a$-cocycle.  Compactness of $H^a$ supplies $\varepsilon>0$ with
$v(z)\geq\varepsilon$ on the entire cochain domain.  Hence
$z^{p^n}\to0$ uniformly.  For $a>0$, apply \Cref{h3:co:strictness} to the
complete coefficient complex, prescribing the open neighborhood of
positive-valued $(a-1)$-cochains.  A sufficiently high Frobenius power
has a primitive $b$ in that neighborhood.  Inverse Frobenius is a
continuous additive automorphism of the completed perfect algebra and
preserves positivity, so $b^{1/p^n}$ is a positive primitive of $z$.
For $a=0$, the finite invariant group is discrete; the invariant elements
$z^{p^n}$ tend to zero and are eventually zero.  Injectivity of Frobenius
then gives $z=0$.

The quotient of $\mathcal O_{B^b}$ by its strictly positive ideal is the
finite product of the residue fields of $B$.  The two valuation-cutoff
sequences are exact on cochains because their quotients are discrete.
Their long exact sequences, the vanishing just proved, and finite
cohomology of finite coefficients establish the last assertions.
\end{proof}

Fix a good $B$ and put
\[
 P=P_B,\qquad M_a=H^a_{\hthreects}(H,P),\qquad
 X_a=H^a_{\hthreects}(H,B),\qquad f_a:M_a\longrightarrow X_a.
\]
The group $M_a$ is compact Hausdorff by
\Cref{h3:co:compact-resolution}.  Since $B/P$ is countable discrete, the
countability observation above and the cochain-exact lattice sequence
give
\begin{equation}\label{h3:co:countable-kernel-cokernel}
 \ker f_a\ \text{and}\ \operatorname{coker}f_a
 \quad\text{are countable.}
\end{equation}

We first prove that every $X_a$ is countable.  If $a$ were the largest
degree where countability failed, the Cartier exact sequences would make
$M_a/\mathcal C M_a$ finite.  The next two lemmas combine this bound with
a finite stable Cartier image to control the Cartier kernel in $X_a$.
Since $\mathcal C\mathrm{Fr}=0$, the Frobenius image would then be
countable; the following kernel estimate makes its kernel countable as
well, a contradiction.  Once every $X_a$ is countable, the comparison
$f_a$ makes each compact $M_a$ finite.  A second use of residue duality
gives finite cohomology of the lattice quotient, and the lattice sequence
then makes $X_a$ finite.

\begin{lemma}[Finite stable Cartier image]\label{h3:co:stable-image}
For every $a$, the subgroup
\[
 I_a=\bigcap_{n\geq0}\mathcal C^nM_a
\]
is finite.
\end{lemma}

\begin{proof}
The residue pairing and \Cref{h3:co:good-stages} give topological
identifications
\begin{equation}\label{h3:co:residue-duals}
 P^\vee=B/\mathcal O_B,\qquad (B/P)^\vee=\mathcal O_B.
\end{equation}
A continuous functional on $P$ uses only finitely many Laurent
coefficients and is represented by a finite principal part in the first
quotient.  An arbitrary functional on the discrete $B/P$ is represented
by a power series, which proves the second identification and its
topology.  The trace transposition in \Cref{h3:co:cartier-trace} and
\eqref{h3:co:compact-duality} identify the dual of the inverse system
$(M_a,\mathcal C)$ with
\[
 \bigl(H^{8-a}_{\hthreects}(H,B/\mathcal O_B),\mathrm{Fr}\bigr).
\]
Its coefficient colimit is
\begin{equation}\label{h3:co:discrete-perfection-quotient}
 \colim_{\mathrm{Fr}}B/\mathcal O_B
 =B^{\mathrm{perf}}/\mathcal O_{B^{\mathrm{perf}}}
 =B^b/\mathcal O_{B^b}.
\end{equation}
At index $n$ the first map sends $f$ to $f^{1/p^n}$.
The second equality follows by approximating any completed element with
integral error.  All these quotients are discrete, so their continuous
cochains commute with the colimit.  By
\Cref{h3:co:completed-positive}, the direct limit of the displayed
cohomology system is finite.  Compact--discrete duality shows that
$\varprojlim_{\mathcal C}M_a$ is finite.

Its zeroth projection has image $I_a$.  To prove surjectivity onto that
intersection, fix $x\in I_a$.  For each $n$ there is a compatible
length-$n$ chain of preimages ending at $x$.  These conditions define
nested nonempty closed subsets of the compact product
$\prod_{n\geq0}M_a$.  Their intersection is an infinite compatible
chain.  Thus $I_a$ is an image of the finite inverse limit.
\end{proof}

\begin{lemma}[A compact Cartier module]\label{h3:co:compact-cartier}
Let $M$ be a compact Hausdorff $\F_p$-vector space with a continuous
endomorphism $C$.  Suppose that $I=\bigcap_n C^nM$ and $M/CM$ are finite.
Then $C(I)=I$, the quotient $Q=M/I$ is a finitely generated
$\F_p[[t]]$-module with $t$ acting as $C$, and both $Q[C^\infty]$ and
$\ker(C:M\to M)$ are finite.
If a $C$-stable subgroup $N\subset M$ has finite forward $C$-orbits
elementwise, then $N$ is finite and
\begin{equation}\label{h3:co:stable-quotient}
 \bigcap_n\hthreeim(C^nM\longrightarrow M/N)
 =\hthreeim(I\longrightarrow M/N).
\end{equation}
\end{lemma}

\begin{proof}
For $x\in I$, the sets $C^{-1}(x)\cap C^nM$ are nested nonempty compact
sets, so they meet.  This proves $C(I)=I$.
The compact images $C^nQ$ have intersection zero and shrink uniformly
to zero: otherwise their intersections with the complement of a fixed
open neighborhood would be nested nonempty compact sets.  Therefore
$\sum_{n\geq0}a_n C^n x$ converges for every formal power series
$\sum a_nt^n$, uniformly in the tails, and defines an
$\F_p[[t]]$-action on $Q$.

Lift a finite basis of $Q/CQ$ to $e_1,\dots,e_s$.  Repeatedly expressing
an element modulo $CQ$ gives
\[
 x=\sum_{i=1}^s\sum_{j=0}^{N-1}a_{ij}C^je_i+C^Nx_N.
\]
The remainder tends uniformly to zero.  Taking the limit expresses $x$
as a sum of $s$ power-series multiples of the $e_i$.  The structure
theorem over the discrete valuation ring $\F_p[[t]]$ now writes $Q$ as
a finite free module plus finitely many modules $\F_p[[t]]/(t^e)$.
Its $C$-power torsion and $C$-kernel are finite.  The kernel on $M$ has
kernel in $I$ and image in the kernel on $Q$, so it is finite too.

The image in $Q$ of an element with finite forward orbit both tends
to zero and belongs to a finite set; it is eventually zero.  Thus the
image of $N$ is in the finite group $Q[C^\infty]$, and its kernel lies
in finite $I$.  Hence $N$ is finite and closed.  A coset meeting every
$C^nM$ meets their intersection, by compactness of its nested
intersections with those images.  This proves
\eqref{h3:co:stable-quotient}.
\end{proof}

We record one algebraic consequence used in applying this lemma.
Suppose $f:M\to X$ commutes with $C$ and has countable kernel and
cokernel.  If $X/CX$ is countable, so is $M/CM$.  Indeed, with
$N=\ker f$, $A=\hthreeim f$, and $E=X/A$, the exact portions
\[
 N/CN\longrightarrow M/CM\longrightarrow A/CA\longrightarrow0,
 \qquad
 \ker(C|E)\longrightarrow A/CA\longrightarrow X/CX
\]
prove the assertion.  If $N$ and $\ker(C|M)$ are finite, then
$\ker(C|A)$ is finite, by its exact sequence with $N/CN$, and
$\ker(C|X)$ is countable since its remaining quotient embeds in $E$.

\begin{lemma}[Frobenius-kernel bound]\label{h3:co:frobenius-kernel}
If $\bigcap_n\hthreeim(\mathcal C^nM_a\to X_a)$ is finite, then
$\ker(\mathrm{Fr}:X_a\to X_a)$ is countable.
\end{lemma}

\begin{proof}
For $a=0$, Frobenius is injective.  For $a>0$, its kernel is the kernel
of the inclusion into the root-stage cohomology under the
$H$-equivariant homeomorphism $B^{1/p}\to B$.  Consequently each kernel
class has a presentation $[\partial b]$, where
\[
 b\in\mathcal B_a
 :=\{b\in C^{a-1}_{\hthreects}(H,B^{1/p}):
               \partial b\in C^a_{\hthreects}(H,B)\}.
\]
Put $w=Db$.  It is a $B^{1/p}$-valued cocycle because $D$ kills $B$.
The quotient map
\[
 \mathcal B_a\longrightarrow
 C^{a-1}_{\hthreects}\bigl(H,B^{1/p}/(B^{1/p})^{++}\bigr),\qquad
 b\longmapsto Db\bmod ++,
\]
has countable target.  Its kernel $\mathcal B_a^+$ consists of
presentations with positive-valued $w$ and has countable index.
It suffices to put the classes presented by this kernel into the finite
intersection in the hypothesis.

Fix such a $b$, and let $i$ range through arbitrarily large positive
integers divisible by two and by $[k:\F_p]$; put $q=p^i$.
For $\theta$ in \Cref{h3:co:good-stages}, define
\[
 b'=\theta^{1/q}w\in C^{a-1}_{\hthreects}(H,B^{1/(pq)}).
\]
By \eqref{h3:co:frobenius-distribution},
$D^q\theta^{1/q}=1$, and $D^q$ is $L^{1/q}$-linear.
As $w\in B^{1/p}\subset L^{1/q}$, this proves $D^qb'=w$.
Taking differentials and using the kernel formula gives
\begin{equation}\label{h3:co:primitive-root-level}
 \partial b'\in C^a_{\hthreects}(H,B^{1/q}).
\end{equation}
Here $b'$ is fixed by $U$, and
$(L^{1/q})^U=B^{1/q}$ by uniqueness of roots.

We must also descend $D^{q-1}b'$.  For $\sigma\in U$ write
$\sigma(D)=D(1+h_\sigma)$, with
$h_\sigma\in D^{p+1}k[[D]]$.  Then
\[
 \sigma(D)^{q-1}-D^{q-1}
 \in D^{q+p}k[[D]].
\]
The whole error kills $b'$, since
$D^{q+p}b'=D^pw=D^{p+1}b=0$ and $D^pb=0$.
Local nilpotence makes each error series finite on the relevant stage.
Thus $D^{q-1}b'$ is $U$-fixed, and
$D(D^{q-1}b'-b)=0$ proves that $D^{q-1}b'-b$ is a $B$-valued
cochain.  For each $q$ all the operators act continuously after one
finite-stage enlargement.  The descended subspaces carry their
valuation subspace topologies, so these descended cochains are
continuous as well.

Let $\alpha_i=(\partial b')^q$.  By
\eqref{h3:co:primitive-root-level}, it is a $B$-valued cocycle.
The descended difference and \Cref{h3:co:cartier-trace} give
\begin{equation}\label{h3:co:kernel-deep-image}
 [\partial b]=[D^{q-1}\partial b']
             =c_i\mathcal C^i[\alpha_i]\quad\text{in }X_a.
\end{equation}
Compactness gives a common $\varepsilon>0$ with $v(w)\geq\varepsilon$.
Because $H$ preserves valuations,
\[
 v(\alpha_i)\geq q\varepsilon+v(\theta).
\]
This eventually exceeds the finitely many thresholds defining $P$.
Hence $\alpha_i$ is a $P$-valued cocycle for all sufficiently large
admissible $i$.  The scalar $c_i$ can be absorbed into the lattice:
$c_i\mathcal C^i(\alpha_i)=\mathcal C^i(c_i^q\alpha_i)$ and
$c_i^qP=P$.  Equation \eqref{h3:co:kernel-deep-image} puts our class in a
cofinal set of the decreasing images $f_a(\mathcal C^iM_a)$, hence
in their intersection.  The subgroup $\mathcal B_a^+$ has finite image
and countable index, so the entire Frobenius kernel is countable.
\end{proof}

\begin{proof}[Proof of \Cref{h3:prop:finite-stage-finiteness}]
Suppose some $X_a$ were uncountable and choose the largest such $a$.
It exists by the degree-eight bound in
\Cref{h3:co:compact-resolution}.  Split \eqref{h3:co:cartier-four-term} at
$V=\hthreeim(d/\eta)=\ker\mathcal C$:
\[
 0\longrightarrow B\xrightarrow{\mathrm{Fr}}B\longrightarrow V
 \longrightarrow0,
 \qquad
 0\longrightarrow V\longrightarrow B\xrightarrow{\mathcal C}B
 \longrightarrow0.
\]
The first sequence places $H^{a+1}_{\hthreects}(H,V)$ between a quotient of
$X_{a+1}$ and a subgroup of $X_{a+2}$, so this group is countable.
The second embeds $X_a/\mathcal C X_a$ into it.  By
\eqref{h3:co:countable-kernel-cokernel} and the observation after
\Cref{h3:co:compact-cartier}, $M_a/\mathcal C M_a$ is countable.
It is compact Hausdorff since $\mathcal C M_a$ is compact and closed.
A countable compact Hausdorff group is finite: Baire applied to its
singleton cover gives an isolated point, and a discrete compact group
is finite.

Apply \Cref{h3:co:compact-cartier} using the finite stable image from
\Cref{h3:co:stable-image}.  To treat $N_a=\ker f_a$, use its connecting
presentation from $H^{a-1}_{\hthreects}(H,B/P)$.  A representing cochain has
finite image, hence a uniform lower bound on the logarithmic orders of
its differential coefficients.  The identity
\[
 \mathcal C^n(f)\eta=\mathrm{Car}^n(f\eta)
\]
and \eqref{h3:co:cartier-laurent} show that for sufficiently large $n$ all
negative indices in that finite image disappear modulo $P$.
Its iterate is therefore a cocycle in the finite, $H$- and
$\mathcal C$-stable module
\[
 P(0)/P,\qquad
 P(0)=\prod_jh_j^{-1}\kappa_j[[s_j]].
\]
Finite cohomology of this module and naturality of the connecting map
show that every element of $N_a$ has finite forward Cartier orbit.
For $a=0$ the kernel is already zero.  The compact-module lemma makes
$N_a$ finite and gives
\[
 \ker(\mathcal C:X_a\to X_a)\text{ countable},\qquad
 \bigcap_n f_a(\mathcal C^nM_a)=f_a(I_a)\text{ finite}.
\]
By \eqref{h3:co:cartier-frobenius-zero}, the Frobenius image in $X_a$ is
countable.  An uncountable group with countable image has uncountable
kernel, contradicting \Cref{h3:co:frobenius-kernel}.

All $X_a$ are therefore countable.  Equation
\eqref{h3:co:countable-kernel-cokernel} makes the compact $M_a$ countable,
hence finite.  Any $H$-stable order lattice is commensurable with $P$;
intersecting the two lattices gives finite quotients.  The associated
long exact sequences prove finite cohomology for every such lattice,
including $\mathcal O_B$.  By the second identification in
\eqref{h3:co:residue-duals} and \eqref{h3:co:compact-duality},
\[
 H^a_{\hthreects}(H,B/P)^\vee
 \cong H^{8-a}_{\hthreects}(H,\mathcal O_B)
\]
is finite.  Thus the quotient cohomology is finite, and the lattice
sequence now proves finiteness of $X_a$ itself.  Other lattice quotients
follow by commensurability.

The homeomorphism $B^{1/p^e}\to B$, $f\mapsto f^{p^e}$, is
$H$-equivariant and $\F_p$-linear and carries order lattices to order
lattices.  It proves the assertion for finite root stages.  Finally,
an $H$-invariant frame identifies a coefficient line at a sufficiently
deep finite stage with $B$ as an $H$-module, and its lattices with order
lattices.  The connected marking supplies such a frame for every
periodicity power; the full norm kernel fixes every determinant twist.
Enlarging the good stage to contain finitely many frame coefficients
proves the stated twisted assertion.
\end{proof}

\subsection{The natural completion comparison and the norm quotient}

Finite cohomology is now known at every cofinal good field and root
stage.  This is the input needed to contract positive valuation
cochains before comparing the algebraic and completed coefficient
complexes.  We then apply the residual norm characters to the resulting
natural comparison.

\begin{lemma}[Positive-order decompletion]\label{h3:co:decompletion}
The inclusions induce a natural quasi-isomorphism
\begin{equation}\label{h3:co:decompletion-map}
 \colim_{U,e}C^\bullet_{\hthreects}(H,B_U^{1/p^e})
 \xrightarrow{\ \sim\ }
 \colim_U C^\bullet_{\hthreects}(H,B_U^b).
\end{equation}
It retains the commuting actions and remains a quasi-isomorphism after
forming continuous families parametrized by any profinite space.
\end{lemma}

\begin{proof}
Fix a good $B$.  The positive order lattice $B^{++}$ has finite
cohomology by \Cref{h3:prop:finite-stage-finiteness}.  A positive-valued
cocycle has a uniform positive valuation bound, so its Frobenius powers
tend uniformly to zero.  \Cref{h3:co:strictness} makes each cohomology class
eventually zero under Frobenius.  Since the cohomology is finite, a
single power in each degree kills every class.

The homeomorphisms $B^{1/p^e}\to B$ identify inclusion of one positive
root stage into the next with Frobenius.  Their cochain colimit is
therefore acyclic.  The positive ideal in $B^b$ is acyclic by
\Cref{h3:co:completed-positive}.  Density gives an isomorphism of discrete
coefficient modules
\[
 B^{\mathrm{perf}}/(B^{\mathrm{perf}})^{++}
 \cong B^b/(B^b)^{++}.
\]
A continuous cochain to either quotient has finite image, whose values
can be lifted at one root stage.  Thus this is also an isomorphism of
the prescribed quotient cochain complexes.  Comparing the two
valuation-cutoff exact sequences proves the quasi-isomorphism for this
$B$.  Passing through the cofinal good stages proves
\eqref{h3:co:decompletion-map}.

Each finite root stage has finite cohomology by
\Cref{h3:prop:finite-stage-finiteness}, and each completed stage has finite
cohomology by \Cref{h3:co:completed-positive}.  Apply
\Cref{h3:co:parameters} to their Polish cochain complexes and pass to the
filtered colimits.  It proves the assertion with any profinite
parameter space.  All comparisons used inclusions and valuation
quotients, so their naturality retains all the indicated actions.
\end{proof}

\begin{lemma}[The full norm weights]\label{h3:lem:full-norm-weights}
On $H^*_{\hthreects}(H,R)$ the full residual norm quotient has the two
characters in \eqref{h3:co:completed-input}.  The distribution target
$R(\chi^{-1})$ has no $H'$-cohomology, whereas
$R\Gamma_{\hthreects}(H',R)$ is $k$ by its constants map.
These assertions hold for every $p\geq5$, without choosing a subgroup
section of $H'\to\mu_{p-1}$.
\end{lemma}

\begin{proof}
\Cref{h3:co:decompletion} transports the actual residual actions in
\eqref{h3:co:completed-input} to $R$.  The quotient $H'/H$ has order
prime to $p$, so its invariants functor is exact; the continuous
Hochschild--Serre complex thus computes $H'$-cohomology by taking those
invariants on $H$-cohomology.  The untwisted characters are $0$ and
$2\epsilon$.  Only the first is trivial modulo $p-1$.
After twisting by $\chi^{-1}$ the exponents are
\[
 -s_{\mathrm{dist}},\qquad 2\epsilon-s_{\mathrm{dist}},\qquad
 \epsilon,s_{\mathrm{dist}}\in\{1,-1\}.
\]
They belong to $\{1,-1,3,-3\}$ and are nonzero modulo $p-1\geq4$.
Both target rows therefore vanish.  In the source the invariant row is
exactly the given constants row.

The distinction from the center can be seen at $p=7$.  A central
$a\in\mu_6\subset P_3$ has norm $a^3$, so pulling a norm-weight-two
character to that central subgroup would make it trivial.  Here the
normalized \'etale generator transforms by the full norm unit $u$;
the source weight is $u^{\pm2}$ and the twisted exponents are
$\pm1,\pm3$, nontrivial modulo six.  The reduced norm surjects on tame
units, as follows from the residue-field norm
$\F_{p^3}^\times\to\F_p^\times$.  Since $H$ acts trivially on the
translation group and by inner automorphisms trivially on its own
cohomology, all residual actions used here factor through the quotient
itself.  Neither the eigenparameter construction nor the invariants
calculation requires it to split inside $P_3$.
\end{proof}

\begin{proof}[Proof of \Cref{h3:thm:coheight-one-comparison}]
The distribution sequence \eqref{h3:co:distribution-sequence} is exact on
the specified continuous cochain complexes.  By
\Cref{h3:lem:full-norm-weights}, its last term has zero $H'$-cohomology,
and its middle term is computed by constants.  Thus its first term is
also computed by the actual constants map, proving
\eqref{h3:co:main-constants}.  The finite-stage $H'$-cohomology is finite
by \Cref{h3:prop:finite-stage-finiteness} and exact tame invariants.
\Cref{h3:co:parameters} proves the parameterized assertion.  Although a
choice of $D$ was used to prove the comparison, the resulting
quasi-isomorphism is the original constants inclusion.  Its naturality
therefore retains the connected-frame action and the remaining
height-three action.
\end{proof}

\subsection{Descent to base fields and base lattices}

The ordinary-stratum argument needs finiteness on the coheight-one
\emph{base lattices}, including all conormal and periodicity twists.
We give descent separately for fields and for lattices; the latter does
not divide by a possibly ramified Galois degree.

\begin{lemma}[Continuous field descent]\label{h3:co:field-descent}
For every profinite space $T$, the natural augmentation is a
quasi-isomorphism
\begin{equation}\label{h3:co:field-descent-map}
 C(T,K)\longrightarrow
 \colim_U C^\bullet_{\hthreects}\bigl(J,C(T,B_U)\bigr).
\end{equation}
The same statement holds for a coefficient line pulled back from $K$.
These comparisons are natural in $T$ and equivariant for the commuting
$P_3$-action.
\end{lemma}

\begin{proof}
For fixed normal open $U$, set $W_U=C(T,B_U)$ with its uniform topology.
The action on it factors through $J/U$.  Compare locally constant
$J$-cochains, viewing $W_U$ as discrete, with continuous $J$-cochains
for its given topology.  A fixed finite-projective
$\F_p[[J]]$-resolution from \Cref{h3:co:compact-resolution} computes both.
Their Hom complexes are the same underlying finite summands of powers
of $W_U$: every module map factors through the finite-dimensional
quotient of a resolution term by the augmentation ideal of $U$, and
is continuous with either topology.  The canonical inclusion of the
two bar-cochain complexes is consequently a quasi-isomorphism.

After taking the colimit over normal open subgroups, every locally
constant cochain, its finite clopen partition, and its finitely many
values are represented at one finite quotient.  Refining $U$ if
necessary gives an equality of complexes
\[
 \colim_U C^\bullet_{\mathrm{lc}}(J,W_U)
 =\colim_U C^\bullet(J/U,W_U).
\]
A finite topological $K$-basis of $B_U$ gives
$W_U=B_U\otimes_K C(T,K)$.  The normal-basis theorem for the finite
Galois \'etale algebra $B_U/K$ makes this a coinduced $J/U$-module.
The finite-group cochain complex of a coinduced module has zero
positive cohomology: its bar contraction is evaluation in the induced
function coordinate.  Its invariants are exactly $C(T,K)$.
Thus the canonical augmentation is a quasi-isomorphism at each finite
quotient.  Exactness of the filtered colimit proves
\eqref{h3:co:field-descent-map}.

For a pulled-back coefficient line, tensor the finite-basis and
normal-basis argument with that line over $K$.  All choices were used
only to prove exactness.  The comparison maps themselves are the
canonical augmentations, hence are natural and retain the commuting
action.  This proof makes no claim about cochains into the metric
union $L$.
\end{proof}

\begin{corollary}[Base lattice finiteness]\label{h3:co:base-lattices}
Let $\mathcal L$ be a coefficient line over $K$ whose pullback to $L$
has an $H$-invariant frame at a finite stage.  Suppose that
$\Lambda\subset\mathcal L$ is a $P_3$-stable order lattice over
$k[[y]]$.  Then
\[
 H^a_{\hthreects}(P_3,\Lambda)
 \quad\text{and}\quad H^a_{\hthreects}(H,\Lambda)
\]
are finite for every $a$.  This applies to every periodicity power,
every finite determinant character, their inverses, and their products
with conormal or canonical-line twists.  In particular each adjacent
$x$-pole strip used in \Cref{h3:sec:ordinary} has finite $P_3$-cohomology.
\end{corollary}

\begin{proof}
Choose a sufficiently deep good normal stage $B=B_U$ containing the
frame and put $Q=J/U$.  Let $\mathcal O_B$ be its integral closure
lattice and set
\[
 M=\mathcal O_B\otimes_{k[[y]]}\Lambda.
\]
The actions of $H$ and $Q$ commute and $M^Q=\Lambda$.
In the $H$-invariant frame, $M$ is an $H$-stable order lattice at the
stage $B$, so \Cref{h3:prop:finite-stage-finiteness} gives finite
$H$-cohomology for it.

For $b>0$, $H^b(Q,M)$ is a finite group.  To see this without averaging,
the finite-group cochain terms are finitely generated modules over
$k[[y]]$.  Their cohomology is finitely generated as well.  Inverting
$y$ commutes with that cohomology and gives zero in positive degrees
by normal-basis descent on the field algebra.  Hence $H^b(Q,M)$ has
finite length over $k[[y]]$, whose residue field is finite.  The finite
group boundaries are compact and closed, so these finite cohomology
groups have their discrete quotient topology and continuous $H$-action.

Use the bounded finite-projective $H$-resolution together with the
finite-$Q$ bar complex on $M$.  Taking $H$-cohomology first gives finite
groups $H^i(Q,H^j_{\hthreects}(H,M))$ in every bidegree, hence finite total
cohomology in every degree.  Taking $Q$-cohomology first gives
\[
 E_2^{a,b}=H^a_{\hthreects}(H,H^b(Q,M)).
\]
All rows with $b>0$ are finite by
\Cref{h3:co:compact-resolution}.  The bottom term
$E_2^{a,0}=H^a_{\hthreects}(H,\Lambda)$ has no outgoing differential and
only finitely many incoming differentials, each with finite image.
Its surviving quotient is finite because the total cohomology is
finite.  Thus it is finite.  These two spectral sequences compute the
same first-quadrant complex: finite projective Hom is exact on the
compact $Q$-cycle and boundary sequences, and the finite $Q$-cochain
functors are finite products.  No strictness assertion about the still
unknown base-lattice cohomology was used.

Finally $P_3/H=\Z_p^\times$ has finite cohomology on finite discrete
modules: take exact invariants for $\mu_{p-1}$ and the two-term
resolution for $1+p\Z_p$.  The continuous Hochschild--Serre sequence
for $H\subset P_3$ now proves finite $P_3$-cohomology of $\Lambda$.
The residual action on its finite $H$-cohomology is continuous, since
the coefficient lattice is compact and the finite-resolution
cohomology has the quotient topology.

The connected marking trivializes all periodicity powers and the full
norm kernel fixes determinant characters.  The canonical-line and
conormal identifications express the other indicated lines as products
of these.  An adjacent $x$-pole strip is a power-series lattice in $y$
with precisely one such conormal and coefficient-line twist.  It
therefore satisfies the hypotheses just proved.
\end{proof}

\begin{remark}[Finite constant fields]\label{h3:co:constant-descent}
Finite enlargement of $k$ does not lose either these finiteness
statements or the natural comparison maps.  Scalar extension of a
coefficient module, of each order lattice, or of each continuous
cochain complex is a finite direct sum in its given topology.
The continuous trace-one contraction of
\Cref{h3:lem:framework-semilinear} makes the associated semilinear Galois
complex acyclic in positive degrees; its invariants are the original
coefficient complex.  This remains true with profinite parameters,
because finite scalar extension commutes with continuous functions
and the contraction consists of finite sums.
For the original extended stabilizer one retains the finite coefficient
Galois action on the Honda stabilizer as well; the equivariant
augmentations give the corresponding semidirect-product descent.
Thus constants, units, line twists, and base lattice finiteness descend
to the original coefficient field.  No division by the constant-field
extension degree is involved.
\end{remark}
